This paper develops an ethics of intimate relations from the Daoist figure of water --- \zh{上善若水}, ``the highest good is like water'' (\emph{Daodejing}, ch.~8). It asks what it is, concretely and between two people, for love to benefit without contending, to seek the lowly place, and to remain unfilled. The guiding decision is to read these water-virtues \emph{structurally} rather than as a counsel of self-cultivation: not as an ascent toward a selfless or quasi-divine condition, but as properties of how value, power, and desire are arranged in a generative relation. The paper accordingly draws its resources less from the contemplative traditions than from political economy, game theory, political philosophy, feminist ethics, Kantian ethics, and psychoanalysis, and treats the contemplative reading as one voice among many rather than as the paper's frame.

A survey of these traditions is organised around the water-virtues themselves --- non-contention, benefiting the ten thousand things, dwelling in the lowly place, formlessness, and non-fullness --- and shows that each virtue is, in different doctrines, both affirmed and opposed. From these oppositions the paper isolates five standing challenges: that non-contention is irrational (game theory), that it is suicidal absent an enforced structure (Hobbes), that dwelling-below is slave morality (Nietzsche), that taking-the-receiver's-shape dissolves the subject (feminism, psychoanalysis), and that formlessness is lawlessness (Kant). The last is answered in place, by reading the water-good as an \emph{imperfect duty} --- lawful yet without fixed form --- using Kant's own resources.

The dialectical core of the paper answers the remaining challenges in two clusters. Against the charges of self-dissolution and slave morality, it argues that the water-good is not formless yielding but \emph{chosen} good: water takes the vessel's shape (utmost softness) yet wears through stone and breaches dykes (utmost hardness), so that a rigid lower bound distinguishes the generative lowly place from mere collapse. Against the charges of irrationality and suicide, it argues --- following the formal companion \citet{wan2026braid} --- that the water-good never demanded the abolition of power, only the shaping of an ineliminable asymmetry so that its surplus is returned rather than retained. A chapter on practice then carries each water-virtue down into the conduct of an actual intimate relation, pairing every form of the good with the vice it most easily becomes, and closing on the limit of all such guidance: that the good cannot be made into a procedure.

Throughout, the formalism of the companion paper is used as a settled foundation and translated, not rebuilt: non-contention as the vanishing of appropriative (gradient) transport in the holonomy that carries surplus; the good and bad lowly place as the solenoidal and gradient parts of the power force; renewal as the avoidance of dynamical solidification; and --- crucially for a philosophy of intimacy --- the non-substitutable particularity of \emph{this} relation as the local ``flesh'' that topological protection, by construction, cannot reach. The water-good is thus presented as the water-phase of the dark virtue (\zh{玄德}, ch.~51): the same non-possessive generativity the companion paper gives a formal body, here given an ethical and practical one, between two people.
